
\noindent
Zur Bestimmung der maximalen Leistung wird die Abhängigkeit der im Lastwiderstand
umgesetzten Leistung vom Lastleitwert $G_L$ untersucht. Ausgangspunkt ist die in
Abschnitt 2.2.1 hergeleitete Leistungsfunktion:
\begin{equation}
\label{eq:PL_GL}
P(G_L) = I_0^2 \frac{G_L}{(G_L + G_I)^2}
\end{equation}

\noindent
Um das Maximum der Leistung zu bestimmen, wird die erste Ableitung der Leistung
nach dem Lastleitwert $G_L$ gebildet:
\begin{equation}
\label{eq:dP_dGL}
P'(G_L) = I_0^2 \frac{(G_L + G_I)^2 - 2G_L(G_L + G_I)}{(G_L + G_I)^4}
\end{equation}

\noindent
Der Zähler kann vereinfacht werden zu:
\begin{equation}
\label{eq:dP_dGL_simpl}
P'(G_L) = I_0^2 \frac{G_I - G_L}{(G_L + G_I)^3}
\end{equation}

\noindent
Das Maximum der Leistung liegt an der Stelle, an der die Ableitung gleich null ist.
Da der Nenner der Ableitung stets positiv ist, bestimmt allein der Zähler die Lage der
Extremstelle:
\begin{equation}
\label{eq:nullstelle_GL}
G_I - G_L = 0
\end{equation}

\noindent
Daraus folgt unmittelbar:
\begin{equation}
\label{eq:GL_eq_GI}
G_L = G_I
\end{equation}

\noindent
Zur Klassifikation der Extremstelle wird zusätzlich die zweite Ableitung der Leistung
nach dem Lastleitwert betrachtet:
\begin{equation}
\label{eq:d2P_dGL2}
P''(G_L) = -\frac{2I_0^2}{(G_L + G_I)^3}
\end{equation}

\noindent
Da der Nenner für $G_L + G_I > 0$ stets positiv ist, ist die zweite Ableitung negativ.
Somit handelt es sich bei der Extremstelle eindeutig um ein Maximum der Leistung.

\noindent
Durch Einsetzen der Bedingung $G_L = G_I$ in Gleichung \eqref{eq:PL_GL} ergibt sich
die maximale Leistung der Stromquelle:
\begin{equation}
\label{eq:P0_stromquelle}
P_0 = \frac{I_0^2}{4G_I}
\end{equation}

\noindent
Mit der Beziehung $G_I = \frac{1}{R_I}$ kann die maximale Leistung auch in Abhän-
gigkeit des Innenwiderstands angegeben werden:
\begin{equation}
\label{eq:P0_stromquelle_R}
P_0 = \frac{I_0^2 R_I}{4}
\end{equation}
